% Unofficial New York University Poster Template
% a fork of https://github.com/anishathalye/gemini
% also refer to https://github.com/k4rtik/uchicago-poster

% Adapted by Phillip Wang, all credits go to Tao Li


\documentclass[final]{beamer}

% ====================
% Packages
% ====================

\usepackage[T1]{fontenc}
%\usepackage[utf8]{luainputenc}
\usepackage{lmodern}
\usepackage[size=custom, width=59.4,height=42, scale=0.7]{beamerposter}
\usetheme{gemini}
\usecolortheme{cam}
\usepackage{graphicx}
\usepackage{booktabs}
\usepackage{tikz}
\usepackage{pgfplots}
\pgfplotsset{compat=1.14}
\usepackage{anyfontsize}

% ====================
% Lengths
% ====================

% If you have N columns, choose \sepwidth and \colwidth such that
% (N+1)*\sepwidth + N*\colwidth = \paperwidth
\newlength{\sepwidth}
\newlength{\colwidth}
\setlength{\sepwidth}{0.025\paperwidth}
\setlength{\colwidth}{0.3\paperwidth}

% Block without a title, used to continue a block from the previous column
\newenvironment{contblock}
  {\begin{beamercolorbox}[colsep*=0ex]{block body}%
    \usebeamerfont{block body}\justifying\setlength{\parskip}{1ex}}
  {\end{beamercolorbox}\vskip0pt\vspace*{2ex}}

\newcommand{\separatorcolumn}{\begin{column}{\sepwidth}\end{column}}

% ====================
% Title
% ====================

\title{Title}

\author{Department of Physics, Durham University}

\institute[shortinst]{Short abstract.}

% ====================
% Footer (optional)
% ====================

\footercontent{
  %\href{https://github.com}{Github: https://github.com} \hfill
  %Capstone Festival, Abu Dhabi \hfill
  %\href{mailto:taoli@nyu.edu}{faiza@nyu.edu}}
}
% (can be left out to remove footer)

% ====================
% Logo (optional)
% ====================

% use this to include logos on the left and/or right side of the header:
% Left: institution
 \logoleft{\includegraphics[height=5cm]{logos/durham-university-logo-transparent.png}}
% Right: funding agencies and other affilations
%\logoright{\includegraphics[height=7cm]{logos/NSF.eps}}
% ====================
% Body
% ====================

\begin{document}



\begin{frame}[t]
\begin{columns}[t]
\separatorcolumn

\begin{column}{\colwidth}

  \begin{block}{Atomic Qubits}

    Modern quantum computing revolves around interactions between two-state quantum systems known as qubits.
    A qubit governed by a time-independent hamiltonian $H_0$ is represented by a superposition state
    $|\psi(t)\rangle = a|0\rangle e^{-i E_0 t / \hbar} + b|1\rangle e^{-i E_1 t / \hbar}$, where $|0\rangle$ and
    $|1\rangle$ are eigenstates of $H_0$ with energies $E_0$ and $E_1$ respectively and $a$ and $b$ are complex
    coefficients. Hamiltonians used
    in quantum computing often contain a time-dependent component $H'(t)$, leading to states with time-dependent
    coefficients $a=a(t)$ and $b=b(t)$. We examine the case of an atomic qubit interacting with a
    plane-polarised electromagnetic (EM) wave with angular frequency $\omega$, with $|0\rangle$ and $|1\rangle$ representing
    the ground state and an excited state of the atom. This interaction perturbs the atom's hamiltonian $H_0$
    by $H'(t)=-\mathbf{d} \cdot \mathbf{\epsilon}(t)$, where $\mathbf{d}$ is the dipole operator and $\mathbf{\epsilon}(t)$ is the
    electric field component of the EM field. By substituting $|\psi(t)\rangle$ with time-dependent $a(t), b(t)$ into the
    time-dependent Schr\"odinger equation with perturbed hamiltonian $H(t)=H_0 + H'(t)$, making the rotating-wave approximation
    and transforming into a rotating
    frame using a change of variables, a set of coupled equations for the rotated $a(t)$ and $b(t)$ is derived. These equations
    take the form of a time-dependent Schr\"odinger equation with hamiltonian
    $H_{int}=\hbar/2
    \begin{pmatrix}
    \Delta & \Omega e^{-i \phi_L} \\
    \Omega e^{i \phi_L} & -\Delta
    \end{pmatrix}$
    in the $|0\rangle$, $|1\rangle$ basis, where $\Delta=\omega-\omega_0$ is the detuning between the EM field and the atomic
    transition's natural frequency $\omega_0=(E_1-E_0)/\hbar$, the Rabi frequency is defined as
    $\Omega=\langle 1|H'(t)|0\rangle / \hbar$, and $\phi_L$ is a constant phase factor in the electric field. As a result, the
    atomic qubit under interaction with the EM field is represented as $|\psi(t)\rangle=\begin{pmatrix}
    a(t) \\
    b(t) \\
    \end{pmatrix}$, with $a(t), b(t)$ governed by $H_{int}$. A single qubit is not useful for quantum computation so we consider
    a system of two atomic qubits interacting with an EM field. On their own, two neutral atoms do not interact with enough
    strength to produce entanglement. A solution is to excite an atom's valence electron to a high-principal number state, such
    as $n>50$, known as a Rydberg state. Such states readily interact with eachother through the van der Waals interaction.
    Denoting the system's four basis states, the possible combinations of each atom being in either its ground state $|0\rangle$
    or a Rydberg state $|r\rangle$,  as $|00\rangle, |0r\rangle, |r0\rangle, |rr\rangle$, the van der Waals interaction
    shifts the energy of the doubly excited state $|rr\rangle$ by $\hbar V\propto R^{-6}$. Extending the single-atom $H_{int}$
    to this two-atom basis and including the van der Waals shift gives a hamiltonian of the form






  \end{block}





\end{column}

\separatorcolumn

\begin{column}{\colwidth}
  $$ H_{int}=
    \begin{bmatrix}
    2\Delta & \Omega e^{-i \phi_L} & \Omega e^{-i \phi_L} & 0 \\
    \Omega e^{i \phi_L} & 0 & 0 & \Omega e^{-i \phi_L} \\
    \Omega e^{i \phi_L} & 0 & 0 & \Omega e^{-i \phi_L} \\
    0 & \Omega e^{i \phi_L} & \Omega e^{i \phi_L} & -2\Delta + 2V
    \end{bmatrix}.
    $$ dfjdskfldsj

  \begin{block}{A block containing an enumerated list}
    $$ H_{int}=
    \begin{bmatrix}
    2\Delta & \Omega e^{-i \phi_L} & \Omega e^{-i \phi_L} & 0 \\
    \Omega e^{i \phi_L} & 0 & 0 & \Omega e^{-i \phi_L} \\
    \Omega e^{i \phi_L} & 0 & 0 & \Omega e^{-i \phi_L} \\
    0 & \Omega e^{i \phi_L} & \Omega e^{i \phi_L} & -2\Delta + 2V
    \end{bmatrix}.
    $$

    Vivamus congue volutpat elit non semper. Praesent molestie nec erat ac
    interdum. In quis suscipit erat. \textbf{Phasellus mauris felis, molestie
    ac pharetra quis}, tempus nec ante. Donec finibus ante vel purus mollis
    fermentum. Sed felis mi, pharetra eget nibh a, feugiat eleifend dolor. Nam
    mollis condimentum purus quis sodales. Nullam eu felis eu nulla eleifend
    bibendum nec eu lorem. Vivamus felis velit, volutpat ut facilisis ac,
    commodo in metus.

    \begin{enumerate}
      \item \textbf{Morbi mauris purus}, egestas at vehicula et, convallis
        accumsan orci. Orci varius natoque penatibus et magnis dis parturient
        montes, nascetur ridiculus mus.
      \item \textbf{Cras vehicula blandit urna ut maximus}. Aliquam blandit nec
        massa ac sollicitudin. Curabitur cursus, metus nec imperdiet bibendum,
        velit lectus faucibus dolor, quis gravida metus mauris gravida turpis.
      \item \textbf{Vestibulum et massa diam}. Phasellus fermentum augue non
        nulla accumsan, non rhoncus lectus condimentum.
    \end{enumerate}

  \end{block}

  \begin{block}{Fusce aliquam magna velit}

    Et rutrum ex euismod vel. Pellentesque ultricies, velit in fermentum
    vestibulum, lectus nisi pretium nibh, sit amet aliquam lectus augue vel
    velit. Suspendisse rhoncus massa porttitor augue feugiat molestie. Sed
    molestie ut orci nec malesuada. Sed ultricies feugiat est fringilla
    posuere.

    \begin{figure}
      \centering
      \begin{tikzpicture}
        \begin{axis}[
            scale only axis,
            no markers,
            domain=0:2*pi,
            samples=100,
            axis lines=center,
            axis line style={-},
            ticks=none]
          \addplot[red] {sin(deg(x))};
          \addplot[blue] {cos(deg(x))};
        \end{axis}
      \end{tikzpicture}
      \caption{Another figure caption.}
    \end{figure}

  \end{block}

  \begin{block}{Nam cursus consequat egestas}

    Nulla eget sem quam. Ut aliquam volutpat nisi vestibulum convallis. Nunc a
    lectus et eros facilisis hendrerit eu non urna. Interdum et malesuada fames
    ac ante \textit{ipsum primis} in faucibus. Etiam sit amet velit eget sem
    euismod tristique. Praesent enim erat, porta vel mattis sed, pharetra sed
    ipsum. Morbi commodo condimentum massa, \textit{tempus venenatis} massa
    hendrerit quis. Maecenas sed porta est. Praesent mollis interdum lectus,
    sit amet sollicitudin risus tincidunt non.

    Etiam sit amet tempus lorem, aliquet condimentum velit. Donec et nibh
    consequat, sagittis ex eget, dictum orci. Etiam quis semper ante. Ut eu
    mauris purus. Proin nec consectetur ligula. Mauris pretium molestie
    ullamcorper. Integer nisi neque, aliquet et odio non, sagittis porta justo.

    \begin{itemize}
      \item \textbf{Sed consequat} id ante vel efficitur. Praesent congue massa
        sed est scelerisque, elementum mollis augue iaculis.
        \begin{itemize}
          \item In sed est finibus, vulputate
            nunc gravida, pulvinar lorem. In maximus nunc dolor, sed auctor eros
            porttitor quis.
          \item Fusce ornare dignissim nisi. Nam sit amet risus vel lacus
            tempor tincidunt eu a arcu.
          \item Donec rhoncus vestibulum erat, quis aliquam leo
            gravida egestas.
        \end{itemize}
      \item \textbf{Sed luctus, elit sit amet} dictum maximus, diam dolor
        faucibus purus, sed lobortis justo erat id turpis.
      \item \textbf{Pellentesque facilisis dolor in leo} bibendum congue.
        Maecenas congue finibus justo, vitae eleifend urna facilisis at.
    \end{itemize}

  \end{block}

\end{column}

\separatorcolumn

\begin{column}{\colwidth}

  \begin{exampleblock}{A highlighted block containing some math}

    A different kind of highlighted block.

    $$
    \int_{-\infty}^{\infty} e^{-x^2}\,dx = \sqrt{\pi}
    $$

    Interdum et malesuada fames $\{1, 4, 9, \ldots\}$ ac ante ipsum primis in
    faucibus. Cras eleifend dolor eu nulla suscipit suscipit. Sed lobortis non
    felis id vulputate.

    \heading{A heading inside a block}

    Praesent consectetur mi $x^2 + y^2$ metus, nec vestibulum justo viverra
    nec. Proin eget nulla pretium, egestas magna aliquam, mollis neque. Vivamus
    dictum $\mathbf{u}^\intercal\mathbf{v}$ sagittis odio, vel porta erat
    congue sed. Maecenas ut dolor quis arcu auctor porttitor.

    \heading{Another heading inside a block}

    Sed augue erat, scelerisque a purus ultricies, placerat porttitor neque.
    Donec $P(y \mid x)$ fermentum consectetur $\nabla_x P(y \mid x)$ sapien
    sagittis egestas. Duis eget leo euismod nunc viverra imperdiet nec id
    justo.

  \end{exampleblock}

  \begin{block}{References}

    Class aptent taciti sociosqu ad litora torquent per conubia nostra, per
    inceptos himenaeos. Phasellus libero enim, gravida sed erat sit amet,
    scelerisque congue diam. Fusce dapibus dui ut augue pulvinar iaculis.

    \begin{table}
      \centering
      \begin{tabular}{l r r c}
        \toprule
        \textbf{First column} & \textbf{Second column} & \textbf{Third column} & \textbf{Fourth} \\
        \midrule
        Foo & 13.37 & 384,394 & $\alpha$ \\
        Bar & 2.17 & 1,392 & $\beta$ \\
        Baz & 3.14 & 83,742 & $\delta$ \\
        Qux & 7.59 & 974 & $\gamma$ \\
        \bottomrule
      \end{tabular}
      \caption{A table caption.}
    \end{table}

    Donec quis posuere ligula. Nunc feugiat elit a mi malesuada consequat. Sed
    imperdiet augue ac nibh aliquet tristique. Aenean eu tortor vulputate,
    eleifend lorem in, dictum urna. Proin auctor ante in augue tincidunt
    tempor. Proin pellentesque vulputate odio, ac gravida nulla posuere
    efficitur. Aenean at velit vel dolor blandit molestie. Mauris laoreet
    commodo quam, non luctus nibh ullamcorper in. Class aptent taciti sociosqu
    ad litora torquent per conubia nostra, per inceptos himenaeos.



  \end{block}

  \begin{block}{References}

    \nocite{*}
    \footnotesize{\bibliographystyle{plain}\bibliography{poster}}

  \end{block}

\end{column}

\separatorcolumn
\end{columns}
\end{frame}

\end{document}
